% !PRINTED 603
\weberpage{603}

If $Q'$ is equal to $A$, then $Q$ is either likewise equal to $A$ or equal to $S$. But if $Q'$ is $=1$, then $Q$ contains no permutation of the first kind besides the unit. If, then, $\varkappa$ and $\lambda$ are two permutations of $Q$, different from one another and different from the unit, then $\varkappa^2$ and $\varkappa\lambda$, being of the first kind, must be $=1$; that is, $\lambda$ must be $=\varkappa$. Thus $Q$ can contain at most one permutation $\varkappa$ different from the identical one. But since $Q$ is to be a normal divisor of $S$, for every permutation $\pi$ of the symmetric group we must have $\pi^{-1}\varkappa\pi = \varkappa$; that is, $\varkappa$ may not change when any permutation whatever is carried out in its cycles. But this is possible only when there are only two figures 1, 2 present at all, and $\varkappa = (1,\,2)$. But then $1, (1,\,2)$ is the whole group $S$.

\begin{center}
\S.\,178.

\gesperrt{Non-metacyclic equations in the field of the rational numbers.}
\end{center}

By the theorem of the preceding section the proof has been given that an equation of degree $n$, when $n$ is greater than 4, can no longer be solved algebraically if the coefficients are regarded as independent variables. But of still greater interest is the question whether there are, in the field of the rational numbers, equations of degree $n$ that are not algebraically solvable. The question can be put somewhat more generally, namely thus: whether there are equations with integral coefficients whose group is the symmetric group, and which therefore, according to a mode of expression explained earlier, have no affect.

If one forms the \gesperrt{Galois} resolvent $G(t)=0$ of degree $\Pi(n)$ of a general equation of degree $n$, $f(x)=0$, with indeterminate coefficients $a$, then $G(t)$ is an integral function of the variables $t$ and $a$ which cannot be decomposed into factors that are again rational functions of $t$ and $a$. If one substitutes for the variables $a$ quantities of any field $\Omega$, and $G(t)$ then becomes reducible in this field, the equation $f(x)=0$ has an affect in the domain of rationality $\Omega$. Our question thus comes down to whether one can put for the variables $a$ such rational numbers in $G(t)$ that $G(t)$ remains irreducible in the field of the rational numbers.
